% ==============================================================
% Section 3: Methodology
% Author: Phúc (LR - LLM Runner)
% File:   paper/sections/03_method.tex
%
% Sources of every number in this section (do not invent values):
%   - data/full_sample.csv, data/pilot_sample.csv
%   - data/full_ground_truth.csv, data/pilot_ground_truth.csv
%   - results/full_api_log.txt, results/pilot_api_log.txt
%   - results/full_llm_output.csv, results/pilot_llm_output.csv
%   - prompt_final.md (v1.0), prompt_v3_candidate.md (v3.0-candidate-s2r-relaxed)
%   - notes.md
%
% NOTE FOR HÙNG (references.bib): this section cites `cohen1960`,
% `landis1977` and `kang2024libro`, which are not in references.bib yet.
% Please add those three entries, otherwise BibTeX prints [?].
%
% NOTE FOR THÊM / QUÂN: this section refers to \ref{sec:results} and
% \ref{sec:threats}. Please use exactly these labels:
%   04_results.tex -> \section{Results}\label{sec:results}
%   06_threats.tex -> \section{Threats to Validity}\label{sec:threats}
%
% NOTE FOR HÙNG/HUY: the random seed used for sampling is not recorded in
% notes.md yet -> see the TODO below before we submit.
% ==============================================================

\section{Methodology}\label{sec:method}

This section describes the experimental design, the dataset, the exact
LLM configuration, and the execution pipeline used to answer our
research question. The protocol was fixed in the approved proposal
before any data was inspected. Every deviation from that protocol is
reported explicitly in Section~\ref{sec:method:amendment}.

% ------------------------------------------------------------------
\subsection{Study Design and Pipeline Overview}
\label{sec:method:design}

We conduct an observational agreement study: an LLM and human
annotators independently label the same bug reports along three
quality dimensions---Observed Behavior (OB), Expected Behavior (EB),
and Steps to Reproduce (S2R)---and we quantify how closely the LLM
reproduces the human consensus. The LLM is therefore treated as a
candidate \emph{quality gate} placed in front of developer triage,
not as a text generator.

The pipeline consists of five stages:

\begin{enumerate}
  \item \textbf{Data collection.} Bug reports are mined from three
        open-source GitHub projects and stored verbatim.
  \item \textbf{Ground truth.} Two annotators label each report using
        a shared rubric; disagreements are resolved into a consensus
        label.
  \item \textbf{LLM labeling.} Each report is sent to the model in a
        single, stateless request that returns a strict JSON object.
  \item \textbf{Parsing and validation.} Responses are parsed; any
        response that is empty or not valid JSON is marked
        \texttt{INVALID} and is never filled in manually.
  \item \textbf{Agreement analysis.} Cohen's Kappa is computed between
        the two annotators (human--human) and between the LLM and the
        developer consensus (LLM--developer).
\end{enumerate}

We ran the pipeline twice: a \emph{pilot} on a subset to validate the
technical setup, and a \emph{full run} on the complete dataset.

% ------------------------------------------------------------------
\subsection{Subject Systems and Dataset}
\label{sec:method:dataset}

We selected three widely used open-source projects that differ in
domain and in reporting culture: \texttt{pandas} (data analysis
library), \texttt{scikit-learn} (machine-learning library), and
\texttt{vscode} (a large IDE with a very high issue volume). This
choice lets us observe whether agreement is stable across a
library-style and an application-style project.

The final dataset contains 250 bug reports; Table~\ref{tab:dataset}
reports its composition. Reports were drawn from issues created
between February 2024 and June 2026, and both open and closed issues
are included (156 closed, 94 open), because closing status is not
part of the construct we measure.

% TODO(Phúc + Hùng): the random seed used for sampling must be recorded
% in notes.md and stated here before submission. Do not guess it.
The pilot set (30 reports) was drawn at random with a fixed seed and
is a strict subset of the full set: all 30 pilot reports were re-labeled
during the full run, so the full run is internally consistent and no
report is labeled under two different configurations in the reported
results.

Only the issue \emph{title} and \emph{body} are given to the model and
to the annotators. Comments, linked pages, commit history, and the
eventual fix are excluded, since a quality gate must decide at
submission time, when none of that information exists yet. Report
bodies average 3{,}668 characters (median 2{,}760).

\begin{table}[t]
\caption{Dataset composition (full run, $N=250$)}
\label{tab:dataset}
\centering
\begin{tabular}{lrrl}
\hline
\textbf{Project} & \textbf{Full} & \textbf{Pilot} & \textbf{Issue creation range} \\
\hline
\texttt{pandas}       &  50 &  5 & 2025-11-05 -- 2026-06-29 \\
\texttt{scikit-learn} &  96 &  9 & 2024-02-08 -- 2026-06-16 \\
\texttt{vscode}       & 104 & 16 & 2026-05-27 -- 2026-06-26 \\
\hline
\textbf{Total}        & \textbf{250} & \textbf{30} & \\
\hline
\end{tabular}
\end{table}

% ------------------------------------------------------------------
\subsection{Annotation Rubric and Ground Truth}
\label{sec:method:groundtruth}

Each of OB, EB, and S2R receives exactly one of four labels:
\texttt{Sufficient} (present and clear enough to act on),
\texttt{Ambiguous} (present but vague), \texttt{Missing} (not stated),
or \texttt{Incorrect} (stated but contradictory or unusable). The three
dimensions are judged independently, and annotators are explicitly
forbidden to reward domain knowledge that is not written in the report
text---the same restriction the model operates under, which keeps the
two conditions comparable.

Annotator~1 labeled every report. Annotator~2 independently labeled a
random subset for inter-annotator agreement: 10 of 30 reports (33\%)
in the pilot and 75 of 250 reports (30\%) in the full run.
Disagreements were resolved by discussion into the consensus columns,
which serve as ground truth for the LLM comparison.

% ------------------------------------------------------------------
\subsection{LLM Configuration}
\label{sec:method:llm}

Table~\ref{tab:config} lists the configuration. We use GPT-4o~mini,
which is a deliberate choice rather than a convenience one: a quality
gate would run on every incoming issue, so cost per report is a design
constraint, and a small model is the realistic deployment candidate.
The same model family has been used in recent bug-report work by
Akyol~et~al.~\cite{akyol2026improbr}.

We set \texttt{temperature}~$=0.0$ to make the labeling as close to
deterministic as the API allows. This differs from generation-oriented
studies, which use sampling to obtain diverse candidates
(e.g., $t=0.7$ in LIBRO~\cite{kang2024libro}); for a classification
gate, output stability is preferable to diversity. Each report is one
independent, stateless request: no conversation history, no examples in
the prompt (zero-shot), and no self-consistency voting.

The system prompt defines OB, EB, and S2R, defines the four labels,
forbids the use of outside knowledge and of inference over absent
information, and requires the reply to be strictly the JSON object
\texttt{\{"OB":\{"label","reason"\},"EB":\{...\},"S2R":\{...\}\}} with
no markdown and no commentary. The user message contains only the issue
title and body. Requiring a machine-parsable JSON schema follows the
prompt-design practice reported by Bo~et~al.~\cite{bo2024chatbr} and
Acharya and Ginde~\cite{acharya2025enhance}, and it is what makes an
\texttt{INVALID} response detectable rather than silently
mis-interpreted.

\begin{table}[t]
\caption{Run configuration (identical for pilot and full run except the prompt version)}
\label{tab:config}
\centering
\begin{tabular}{ll}
\hline
\textbf{Parameter} & \textbf{Value} \\
\hline
Model (requested)   & \texttt{gpt-4o-mini} \\
Model (returned)    & \texttt{gpt-4o-mini-2024-07-18} \\
Temperature         & \texttt{0.0} \\
Prompting strategy  & zero-shot, single call per report \\
Output format       & strict JSON, no markdown \\
Prompt (pilot)      & \texttt{v1.0} \\
Prompt (full run)   & \texttt{v3.0-candidate-s2r-relaxed} \\
Input given to model& issue title + body only \\
\hline
\end{tabular}
\end{table}

% ------------------------------------------------------------------
\subsection{Prompt Amendment Between Pilot and Full Run}
\label{sec:method:amendment}

The pilot used prompt \texttt{v1.0}. Pilot agreement was acceptable for
OB and EB but low for S2R, and manual inspection of the disagreements
showed a systematic cause: \texttt{v1.0} led the model to expect
explicit, enumerated steps, so reports whose reproduction path was a
runnable code snippet or a single command---the dominant style in
\texttt{pandas} and \texttt{scikit-learn}---were labeled
\texttt{Missing}, whereas the annotators accepted them as actionable.

The full run therefore used prompt \texttt{v3.0-candidate-s2r-relaxed},
which changes two things relative to \texttt{v1.0}:

\begin{itemize}
  \item \textbf{S2R relaxed:} a concrete code snippet, command block, or
        short workflow that identifies the triggering action counts as
        \texttt{Sufficient} even without numbered steps; \texttt{Missing}
        is reserved for reports with no actionable reproduction path at all.
  \item \textbf{EB tightened:} EB is \texttt{Sufficient} only when the
        report explicitly states what should have happened; an
        expectation that is merely guessable from context is not enough.
\end{itemize}

We report this change explicitly because it is a deviation from the
frozen-prompt protocol, and we make three points about its scope.
First, the change is an alignment of the prompt with the annotation
rubric, not a change of the rubric, the metric, the threshold, or the
model. Second, it was decided from pilot \emph{error patterns} and
applied to all 250 reports uniformly, including the 30 pilot reports,
so the reported full-run results come from a single configuration.
Third, because pilot and full run differ in both prompt version and
sample size, pilot and full numbers must not be read as a controlled
comparison of prompt versions; we treat this as a threat to validity
(Section~\ref{sec:threats}).

% ------------------------------------------------------------------
\subsection{Execution Pipeline and Error Handling}
\label{sec:method:pipeline}

Because an API-based pipeline can fail in ways that quietly corrupt a
dataset, error handling was defined before the run rather than
improvised during it:

\begin{itemize}
  \item \textbf{Rate limiting.} Requests are retried up to five times
        with exponential backoff and jitter ($2^{n}$~seconds plus a
        random offset). The jitter avoids synchronised retries; the cap
        avoids masking a persistent outage as slowness.
  \item \textbf{Other errors.} Any non-rate-limit exception is raised
        instead of retried, so that a genuine defect surfaces
        immediately rather than being retried into silence.
  \item \textbf{Empty or unparsable responses.} The row is marked
        \texttt{INVALID} in \texttt{raw\_json\_status}. Labels are never
        filled in by hand and invalid rows are never deleted; they are
        counted separately so that the invalid rate is reported next to
        the metric.
  \item \textbf{Logging.} Every call appends one line to the API log
        recording timestamp, repository, issue id, the model string
        actually returned, success/failure, parse status, token usage,
        cost, and any error message. The returned model string is logged
        (not assumed) so that a silent model rollover is detectable.
\end{itemize}

In practice, all 30 pilot calls and all 250 full-run calls succeeded
and returned parsable JSON: 0/250 \texttt{INVALID}, no retries
triggered, and no format change during the run. One anomaly was
observed: 2 of the 250 responses reported the model alias
\texttt{gpt-4o-mini} instead of the pinned snapshot
\texttt{gpt-4o-mini-2024-07-18}; both parsed correctly and were kept,
and we note them because exact snapshot pinning is a reproducibility
concern for API-based studies.

% ------------------------------------------------------------------
\subsection{Measurement Procedure}
\label{sec:method:metrics}

Agreement is measured with Cohen's Kappa~\cite{cohen1960}, which
corrects for the agreement expected by chance. This correction matters
here: the label distribution is heavily skewed toward
\texttt{Sufficient}, so raw agreement alone would overstate the
model's ability, and we therefore report raw agreement only alongside
Kappa, never instead of it. Kappa is also the measure used by recent
bug-report studies that compare human raters, which makes our numbers
comparable to theirs~\cite{mahmud2025astrobr,akyol2026improbr}. We
interpret magnitudes following Landis and Koch~\cite{landis1977}.

We compute Kappa (i) between the two annotators, to establish that the
rubric itself is reliable, and (ii) between the LLM and the developer
consensus, which is the quantity the research question asks about.
Both are computed overall and per dimension (OB, EB, S2R), since a
single overall value can hide one weak dimension. The pass threshold
fixed in the proposal is $\kappa \geq 0.70$. Results are reported in
Section~\ref{sec:results}.

% ------------------------------------------------------------------
\subsection{Cost, Runtime, and Reproducibility}
\label{sec:method:repro}

The pilot (30 calls, 2026-06-29) consumed 50{,}174 prompt tokens and
3{,}545 completion tokens for approximately US\$0.010 in under two
minutes. The full run (250 calls, 2026-07-06) consumed 449{,}688 prompt
tokens and 30{,}894 completion tokens for approximately US\$0.086 in
about 15 minutes, i.e.\ roughly US\$0.0003 per bug report. This
supports the deployment premise: screening an entire issue backlog with
this gate costs cents, so the binding constraint is accuracy, not cost.

All artifacts required to reproduce the study are in the project
repository: the sampled datasets, the frozen prompts, the parsed model
outputs, and the per-call API logs. No API key is stored in the
repository; credentials are read from the environment at runtime.